\chapter{参考文献与延伸阅读}
\begin{enumerate}[label={[\arabic*]}]
\item 茆诗松、程依明、濮晓龙、倪葎，\textit{概率论与数理统计教程}，第四版，高等教育出版社，2026。正文简称“茆诗松”，按章节定位。
\item 李东风，\textit{概率统计讲义}，2016年秋季课程讲义；对应何书元《概率论与数理统计》第一版。正文简称“李东风”。
\item Hao Wu，\textit{Introduction to Probability Theory}，课程讲义。正文简称“吴昊”；可测性、收敛与条件期望按讲义节号定位。
\item Athanasios Papoulis、S. Unnikrishna Pillai，\textit{概率、随机变量与随机过程}，第四版，保铮、冯大政、水鹏朗译，西安交通大学出版社，2012。正文简称“Papoulis--Pillai”。
\item \textit{统计学知识笔记}，数学工具、抽样推断与线性模型章节。正文简称“统计笔记”，按章节定位。
\end{enumerate}
相关标准理论的适用条件与推导均在相应章节给出。
可沿“概率空间—积分—收敛—条件期望”深入概率基础，沿“充分性—完备性—无偏最优—似然比”深入数理统计。
本册题目采用基础、考研风格、竞赛风格和测度论选学标签，不附未经核实的年份、学校或比赛来源。
